
\subsection{Motivation and Problem Statement}
Mobile robots are progressively being deployed beyond controlled indoor spaces to \emph{unstructured} environments such as construction sites, farms, and disaster areas.  In these settings, robots must operate side-by-side with human workers, heavy machinery, and uneven terrain composed of gravel, mud, and loose rubble.  Achieving \emph{socially compliant navigation} such as moving efficiently toward a goal while avoiding unsafe incursions into human work areas, therefore requires reasoning about \emph{both} human intent \emph{and} terrain affordances in real time.  Formally, let $\gamma(t)$ denote a robot trajectory from start $s_{0}$ to goal $g$, and let $\mathcal{H}(t)$ represent the set of predicted human occupancies over the same horizon.  The objective is to find
\begin{equation}
\gamma^{*} \;=\; \argmin_{\gamma}
\;\mathcal{C}_{\text{terrain}}(\gamma)
+\lambda_{1}\,\mathcal{C}_{\text{social}}\bigl(\gamma,\mathcal{H}\bigr)
+\lambda_{2}\,\mathcal{C}_{\text{goal}}(\gamma),
\label{eq:objective}
\end{equation}
subject to the robot’s kinodynamic constraints, where $\mathcal{C}_{\text{terrain}}$ penalises traversing non-drivable ground, $\mathcal{C}_{\text{social}}$ penalises projected interference with humans, and $\mathcal{C}_{\text{goal}}$ rewards task completion.  Solving~\eqref{eq:objective} online in a partially observed, rapidly changing environment remains a fundamental challenge.

\subsection{Limitations of Prior Work}
Indoor social-navigation planners ranging from velocity-obstacle variants~\cite{van2011reciprocal} to deep reinforcement learning (DRL) approaches~\cite{everett2021collision}, generally assume flat terrain and consistent pedestrian motion.  When transferred outdoors, most methods ignore slope, soil cohesion, and wheel–ground interaction.  Conversely, off-road navigation research has produced terrain classifiers~\cite{brooks2023srpa}, height-map planners~\cite{rankin2024vtr}, and RL agents~\cite{replace} that excel at negotiating surface hazards, yet treat humans as static obstacles or omit them entirely.  Human trajectory-forecasting literature offers sophisticated spatiotemporal models~\cite{kosaraju2020social} but seldom integrates their predictions into the feedback loop of robot control.  As a result, no existing methods that we are aware of jointly optimises for socially aware navigation and terrain traversability.

%\subsection{Key Insight and Proposed Approach}
%Human workers implicitly reveal safe, task-relevant regions through their walking patterns: they avoid slippery mud, prefer compacted soil over rocks, and choose ramps over ladders.  We leverage this observation by \emph{learning} a \textbf{human traversability prior} with a lightweight spatiotemporal attention network.  Given a short $T$-frame RGB sequence, the network outputs a heat-map that peaks on areas humans are likely to walk on within a \SI{2}{s} horizon.  This soft prior is fused—with depth and grayscale appearance—into a three-channel observation for a model-based Soft Actor–Critic (SAC) agent that also learns a latent world model of the robot’s own dynamics.  The resulting hybrid pipeline (Fig.~\ref{fig:pipeline}) plans socially compliant, terrain-aware trajectories at \SI{10}{Hz}, without requiring hand-crafted costmaps or explicit semantic labels at test time.

\subsection{Contributions}
This paper makes three primary contributions:
\begin{enumerate}
  \item A soft real-time social naviagtion in unstructured environments solution.
  \item A compact spatiotemporal attention architecture for future state predicter that can be trained on any YOLO object classes.  
  \item An extensive empirical study across \num{3} unstructured scenes with dynamic humans
\end{enumerate}

\subsection{Paper Organisation}
Section~\ref{sec:related} surveys related work on social navigation and unstructured terrain traversability estimation.  Section~\ref{sec:methods} details our methodology, covering dataset preparation, the attention network, SAC formulation, and system integration.  Experimental protocols and quantitative results appear in Section~\ref{sec:results}.  Section~\ref{sec:discussion} analyses limitations and future directions, and Section~\ref{sec:conclusion} concludes.
